\scanpage{291}

This \deleted{\unsure{posited}} \inserted{\unsure{granted}} we have \inserted{defined} a functor
\[
(\pi',\Gamma,I^{\prime !})\longmapsto(\pi,\Gamma,I^{\prime !})
\]
going from the category of outer $\Gamma$-groups with loops
\inserted{$\pi'$} equipped with a subset $I^{\prime !}$ of $I'=I(\pi')$ stable
\inserted{\deleted{on which $\Gamma$ operates} such that $\unsure{s'}\in I'\Rightarrow\Gamma_{\unsure{s'}}\ne(1)$,}
under $\Gamma$, to the category of outer $\Gamma$-groups with loops, equipped with a subset $I^{\prime !}$ of
$\Phi(\pi,\Gamma)$. The conjecture is that the functor
(itself somewhat conjectural, unless we restrict to realizable source situations, whence target situations also realizable) is an equivalence of categories, i.e. $1^\circ)$ fully faithful and $2^\circ)$ \unsure{essentially} surjective. \deleted{The point}

The point $2^\circ)$ is clear, when we restrict on both sides to realizable situations. But even in this case, the point $1^\circ)$ is not proved.
\deleted{\unsure{It may be that}} \inserted{I think} that
\deleted{\unsure{the properties} \illegible{2} \unsure{on the} \illegible{2} \unsure{faithfulness}}
\inserted{\uline{if} $I'$ \uline{is reduced to} \uline{one element}, then the}
\deleted{\[
\operatorname{Autext}(\pi',I^{\prime !})\longrightarrow
\operatorname{Autext}(\unsure{\pi)\Gamma|\pi,\Gamma},\unsure{I^{\prime !}|I''}
\]}
properties of the ``drilling a hole'' functor will allow us to establish it easily. We also see that to establish it, we are

\scanpage{292}
\mspage{147}

reduced to establishing the bijectivity of the $(*)$, and the \unsure{theorem} of equivalence in the case where $\Gamma$ is transitive on $I^{\prime !}$.

To summarize, the Program of attack on the Problem of realizability of an outer $\Gamma$-group with loops would be the following

I) Establish the bijectivity of the $(*)$ in the general case

(II) \unsure{Establish} the full faithfulness of the preceding functor (by \unsure{reducing} if necessary to the case where $I^{\prime !}$ is \uline{one} orbit (singular) of $\Gamma$ in $I'=I(\pi')$.

(III) In the case where $\Phi(\pi,\Gamma)=\varnothing$ i.e. the extension $\unsure{\mathcal E}$ has no element of finite order $\ne1$, prove that $\unsure{\mathcal E}$, equipped \inserted{with the set} of $\unsure{\mathcal E}$-conjugacy classes of the centralizers of the $L_i$ $(i\in I(\pi))$ in $\unsure{\mathcal E}$, is a group with loops.

Moreover, to have the indispensable preliminary properties of the sets $\Phi(\pi,\Gamma)$, we would have to begin by carrying out this program for the

\scanpage{293}

cyclic groups (\unsure{operating} in a direct way), and proceed in this case by dévissage.

If the complete classification theorem (as an equivalence of the 2-categories \inserted{--of operations} that are topological and of \unsure{isotopic|homological} operations --) were true, points I, II, III of this program \unsure{should} be so too, and it would therefore be a good program of approach. Points I and II should be valid, as soon as the 2-functor between 2-categories were 2-faithful (not necessarily \deleted{not} 3-faithful), \deleted{and the necessary condition} \inserted{\unsure{at least for}} realizable situations. We would thus say that an \inserted{\unsure{outer}} action of $\Gamma$ on a $\pi'$, equipped with an $I^{\prime !}$, is \uline{admissible}, if \deleted{every}
\[
\operatorname{Aut\,ext}(\pi',\Gamma,I^{\prime !})\xrightarrow{\sim}\operatorname{Aut}(\pi,\Gamma,I^{\prime !})
\]
and \unsure{if} moreover every \deleted{\unsure{good}} other action of $\Gamma$ on $\pi'$, \deleted{inducing the same action on $I^{\prime !}$ and} the same action on $\pi$ and the same map $I^{\prime !}\unsure{\longrightarrow|{-}}\Phi(\pi,\Gamma)$, is conjugate to the original action

\scanpage{294}
\mspage{148}

This posited, the \unsure{necessary} and sufficient condition for \unsure{realizability} of an outer \inserted{faithful} action of $\Gamma$ on a $\pi$ would then be that this action \unsure{lifts} \inserted{(by ``\unsure{drilling} holes'' \deleted{\unsure{for $I^{\prime !}$}} \inserted{for $I^{\prime !}=\Phi(\pi,\Gamma)$})} to an action \inserted{\uline{admissible},} of $\Gamma$ on a $(\pi',I^{\prime !})$, \deleted{\illegible{1}} (by the very definition, this \unsure{lifting} would be \deleted{that the corresponding extension of $\Gamma$ \unsure{by} \unsure{some} $\pi'$,} then \unsure{essentially} unique) and that moreover the corresponding extension $\unsure{\mathcal E}'$ of $\Gamma$ by $\pi'$ be a group with loops.

25. \uline{Special} Teichmüller groups.

Let us return to the notation of number 19, We are going to define a subgroup $SA_{g,\nu}^{\unsure{\deleted{!}|!}}$ of $A_{g,\nu}^{!}$
\[
SA_{g,\nu}^{\unsure{\deleted{!}|!}}\unsure{\{} =
\{u\in A_g\mid u\text{ induces the identity on a neighborhood of }S_{g,\nu}\}
\]
\inserted{(set of \unsure{automorphisms} of $U_{g,\nu}$ which)}
i.e. \deleted{$u$} \unsure{induce|induces} the identity \unsure{in|in} the complement of a compact set.

Unlike the other subgroups of $A_g$ considered so far, this one is not a closed subgroup. \uline{N.B.} In the case where we work on compact surfaces

\scanpage{295}

with boundary instead of ``\unsure{punctured|punctured}'' compact surfaces (without boundary), it would be appropriate to take the group of automorphisms which \unsure{induce} the identity on the boundary.

$SA_{g,\nu}$ is an \unsure{invariant} subgroup of $A_{g,\nu}$, the \unsure{quotient} $A_{g,\nu}/SA_{g,\nu}$ being \unsure{isomorphic} to the group of \uline{germs} of \unsure{automorphisms} of $X_g$ in a neighborhood of $\unsure{S_{g,\nu}|SA_{g,\nu}}$, or again the group of germs \inserted{at infinity} of \unsure{automorphisms} of $\unsure{U_{\nu,g}|U_{g,\nu}}$ (in the complements of compact sets\ldots)

\deleted{Let us consider} \inserted{We will have} obviously, since
\[
SA_{g,\nu}\subset A_{g,\nu}
\]
\[
(SA_{g,\nu})^{0}\subset \unsure{SA_{g,\nu}^{0}|A_{g,\nu}^{0}},
\]
\unsure{let us put}
\[
S\Gamma_{g,\nu}^{!}=SA_{g,\nu}/SA_{g,\nu}^{0},
\]
we will have a canonical homomorphism
\[
S\Gamma_{g,\nu}^{!}\longrightarrow\Gamma_{g,\nu}^{!+}.
\]
which we are going to interpret algebraically, in terms of the interpretation of $\Gamma_{g,\nu}^{!+}$ \inserted{$=\Gamma^{!}$} as

\scanpage{296}
\mspage{149}

the group of outer automorphisms of the group with loops $\pi_{g,\nu}$, \unsure{inducing} the identity on $S_{g,\nu}=I(\pi_{g,\nu})$.

\deleted{For every \illegible{1}} \deleted{Let} For $i\in I=S_{g,\nu}$, let us choose an \deleted{$\unsure{\ell}$} $L_i$ of class $i$ in $\pi_{g,\nu}$ --- which amounts to giving oneself a ``point'' \unsure{of the|of} $\unsure{B}_{D_i^{*}}$ \inserted{($\simeq$ a universal covering $\widetilde D_i^{*}$)} \deleted{\unsure{on}} and an isomorphism \unsure{between} its image \unsure{in} $\unsure{B}_{U=U_{g,\nu}}$ with the \deleted{\unsure{element}} reference ``point'' $s=s_{g,\nu}$, which served to define $\pi_{g,\nu}$ as $\operatorname{Aut}_{\unsure{B}_U}(s_{\deleted{\unsure{g,\nu}}})\simeq\pi_1(\unsure{B}_U,s)$.

This said, if $u$ is an automorphism of $\pi=\pi_{g,\nu}$, the fact that it respects (\unsure{strictly}, \inserted{by inducing the identity on $S_{g,\nu}$)} the structure with loops, is expressed by the existence of a family of elements $g_i\in\unsure{\pi}$, \unsure{such} that
\[
u(\ell)=\operatorname{int}g_i^{-1}\bigl(\deleted{\unsure{\chi(u)}}\ell^{\deleted{\unsure{\chi(u)}}\inserted{\alpha}}\bigr)
\qquad(i\in I,\unsure{\forall}\ell\in L_i)
\]
\margin{where $\alpha=\chi(u)$}
--- \unsure{which|which} $g_i$ are determined by $u$ \deleted{\illegible{1}} modulo right multiplication by a $\lambda_i\in L_i$. If we likewise have $v,(h_i)$, then \deleted{on $L_i$} For $\ell\in L_i$ \unsure{we have}
\[
\begin{aligned}
(vu)(\ell)
&=v\bigl(\underbrace{u(\ell)}_{\inserted{\operatorname{int}(g_i^{-1})\ell^{\alpha}}}\bigr)
=v\bigl(\operatorname{int}(g_i^{-1})\bigr)v(\ell)^{\alpha}\\
&=\operatorname{int}\bigl(v(g_i^{-1})\bigr)\operatorname{int}\bigl(\unsure{h_i^{-1}}\bigr)\cdot(\ell^{\alpha\beta})\\
&=\operatorname{int}\bigl(v(g_i^{-1})\cdot h_i^{-1}\bigr)\cdot\ell^{\alpha\beta}
=\operatorname{int}\bigl((h_i\cdot v(g_i))^{-1}\bigr)\ell^{\alpha\beta}
\end{aligned}
\]
\margin{N.B. $\alpha=\chi(u)$\quad $\beta=\chi(v)$}
thus $vu$ is compatible with the system of \deleted{$\unsure{u|v}(g_i)$} $h_i v(g_i)$ \unsure{Let us put}

\scanpage{297}

\[
(v,(h_i))\cdot(u,(g_i))=(vu,h_i v(g_i)\deleted{\unsure{h_i}}),
\]
\margin{$S$}

We then find a group structure on the set \inserted{$S\unsure{\mathcal E}^{!}$} of $(u,(g_i))$, a subgroup of the semidirect product of $\unsure{\mathcal E}^{!}=\mathrm{Autext}_{\mathrm{lac}}(\pi,I)$ \inserted{identity on $I$} by $\pi^I$, on which $\unsure{\mathcal E}^{!}$ operates in the obvious way. The natural homomorphism
\[
S\unsure{\mathcal E}^{!}\longrightarrow\unsure{\mathcal E}^{!}
\]
is surjective, and its kernel is canonically isomorphic to $\prod_{i\in I} L_i\simeq \unsure{T|\mathbb T}^I$, where $\unsure{T|\mathbb T}$ is the orientation module.
\deleted{I would like to define a canonical \unsure{homomorphism}
\[
S\Gamma_{g,\nu}^{!}\longrightarrow S^{\prime}\Gamma^{!}
\]
compatible with the \unsure{homomorphisms} of the $\Gamma^{!}$, which defines}
\inserted{whence} an extension structure \deleted{(\unsure{central})}
\margin{where $\unsure{\mathcal E}$ operates on $\unsure{T|\mathbb T}^I$ via its action on $\unsure{T|\mathbb T}$}
\[
1\longrightarrow \unsure{T|\mathbb T}^I\longrightarrow S\unsure{\mathcal E}^{!}\longrightarrow\unsure{\mathcal E}^{!}-1
\]
\deleted{
\[
1\unsure{\longrightarrow}\unsure{T|\mathbb T}_g^{S_{g,\nu}}\unsure{\longrightarrow}S\Gamma_{g,\nu}^{!}\unsure{\longrightarrow}\Gamma_{g,\nu}^{!}-1
\]
\unsure{on} $S\Gamma_{g,\nu}^{!}$ (where $\unsure{T|\mathbb T}_g\,(\unsure{\simeq}\unsure{\mathbb Z})$ is the orientation module). We must therefore define

Note that the kernel of $S\Gamma_{g,\nu}^{!}\longrightarrow\Gamma_{g,\nu}^{!}$ is identified with
$SA_{g,\nu}^{\unsure{\deleted{!}|!}}\cap A_{g,\nu}^{0}/(SA_{g,\nu}^{\unsure{\deleted{!}|!}})^{0}$, and
}

\scanpage{298}
\mspage{150}

\deleted{
\[
SA_{g,\nu}\longrightarrow S^{\prime}\Gamma^{!}
\]
using the choice, for $\unsure{\forall|\mathrm{tt}}\ i\inserted{\in I=S_{g,\nu}}$ of a universal covering $\widetilde D_i^{*}$ of a neighborhood $D_i^{*}$ of $s_i$ (\unsure{which we} will always take: a punctured disk) and of an \unsure{isomorphism} between the universal covering $\unsure{\widetilde U(i)|U(i)}$ of $U=U_{g,\nu}$ which it defines, and the reference \unsure{universal} covering $\widetilde U$ (pointed over $s=s_{g,\nu}$). This \unsure{automorphism}
}

I would like to \unsure{interpret} this extension as the inverse image of a \inserted{canonical} extension of $\Gamma^{!}$ by $\unsure{T|\mathbb T}^I$ (canonical in any case, once the $\widetilde D_i^{*}$ have been chosen). For this, note that if $\unsure{\ddot u|\dot u}$ is an \unsure{autoequivalence|autoequivalence} \deleted{\unsure{in}} of the situation $\unsure{B}_{D^{*}}\longrightarrow\unsure{B}_U$ \inserted{inducing the identity on $I$} we can consider the $\dot u(\widetilde D_i^{*})$, and we can consider the $\widetilde{\dot u}$ over $\dot u$, namely the systems $(\dot u,\gamma_i)$, where for \unsure{every|every} $i\in I$, $\gamma_i:\widetilde D_i^{*}\longrightarrow\dot u(\widetilde D_i^{*})$ [\unsure{determined} \unsure{modulo|by} \unsure{element|element} of $\unsure{T|\mathbb T}$]. \deleted{\unsure{The set of}} The group of isomorphism classes of \unsure{automorphisms|automorphisms} of $(\unsure{B}_{D^{*}}\unsure{\longrightarrow}\unsure{B}_U,(\widetilde D_i^{*}))$, or if we prefer, of isomorphism classes of \unsure{automorphisms|automorphisms} of
\[
(\unsure{B}_I\longrightarrow\unsure{B}_{D^{*}}\unsure{\longrightarrow}\unsure{B}_U)
\]

\scanpage{299}

inducing the identity on $I$, is therefore an extension of $\Gamma^{!}$ by $\unsure{T|\mathbb T}^I$, say $S\Gamma^{!}$. We are going to define a homomorphism of exact sequences
\diagram{\begin{tikzcd}[ampersand replacement=\&,column sep=2em,row sep=2.2em]
1 \arrow[r,no head,to path={(\tikztostart.east) -- ([xshift=-2pt]\tikztotarget.west) -- ++(0,-0.07) \tikztonodes}]
\& T^I \arrow[r,no head] \arrow[d,equal]
\& S\unsure{\mathcal E}^{!}
 \arrow[r,no head,to path={(\tikztostart.east) -- ([xshift=-2pt]\tikztotarget.west) -- ++(0,-0.07) \tikztonodes}]
 \arrow[d,no head]
\& \deleted{S}\unsure{\mathcal E}^{!}
 \arrow[r,no head,to path={(\tikztostart.east) -- ([xshift=-2pt]\tikztotarget.west) -- ++(0,-0.10) \tikztonodes}]
 \arrow[d]
\& 1
\\
\deleted{\unsure{\illegible{1}\,\Gamma}}
 \arrow[r,no head,to path={(\tikztostart.east) -- ([xshift=-2pt]\tikztotarget.west) -- ++(-0.03,-0.06) \tikztonodes}]
\& T^I
 \arrow[r,no head,to path={(\tikztostart.east) -- ([xshift=-2pt]\tikztotarget.west) -- ++(-0.03,-0.06) \tikztonodes}]
\& S\Gamma^{!} \arrow[r,no head]
\& \Gamma^{!} \arrow[r,no head]
\& 1
\end{tikzcd}}
which proves that $S\unsure{\mathcal E|\varepsilon}^{!}$ is indeed the inverse image of \deleted{\illegible{1}} the extension \inserted{$S\Gamma^{!}$} of $\Gamma^{!}$ by $\unsure{T|\mathbb T}^I$.

For this, we note that $\deleted{S}\unsure{\mathcal E|\varepsilon}^{!}$ is the group of isomorphism classes \unsure{of autoequivalences|of autoequivalences} of
\diagram{\begin{tikzcd}[ampersand replacement=\&]
\unsure{B}_{D^*} \arrow[d] \& {} \\
\unsure{B}_U \& \text{point}s \arrow[l]
\end{tikzcd}}
and $S\unsure{\mathcal E|\varepsilon}^{!}$ that of isomorphism classes of \deleted{\illegible{1}} autoequivalences of
\diagram{\begin{tikzcd}[ampersand replacement=\&,column sep=2.4em,row sep=1.7em]
B_{D^*} \arrow[dd]
\& {} \& B_I \arrow[ll] \arrow[dd]
\\
{} \& {} \& {}
\\
B_U
\& {} \arrow[ul,bend right=20,shorten <=5pt,shorten >=5pt,"{\begin{gathered}\text{isomorphism}\\[-2pt]\text{of}\\[-2pt]\text{commutativity}\end{gathered}}"'{xshift=1.3em,yshift=-1.8em}]
\& \text{point}\,s \arrow[ll]
\end{tikzcd}}
which maps naturally to that of isomorphism classes of autoequivalences of the \unsure{2-diagrams}
\diagram{\begin{tikzcd}[ampersand replacement=\&]
\unsure{B}_I \arrow[r] \& \unsure{B}_{D^*} \arrow[r] \& \unsure{B}_U
\end{tikzcd}}.

\scanpage{300}
\mspage{151}

In fact, in all this there was no reason at all to restrict ourselves to the \deleted{\unsure{outer}} \inserted{\unsure{automorphisms}} of $\pi$ (outer or total) inducing the identity on $I$. In any case we find extensions $S\unsure{\mathcal E|\varepsilon}$ of $\unsure{\mathcal E|\varepsilon}$ by $\unsure{T|\mathbb T}^I$, $S\Gamma$ of $\Gamma$ by $\unsure{T|\mathbb T}^I$, and a homomorphism of extensions
\diagram{\begin{tikzcd}[ampersand replacement=\&]
{} \& {} \& 1 \arrow[d] \& 1 \arrow[d] \& {} \\
{} \& {} \& \pi \arrow[r,"\unsure{\sim}"] \arrow[d] \& \pi \& {} \\
1 \arrow[r] \& \unsure{T|\mathbb T}^I \arrow[r] \arrow[d,equal] \& S\unsure{\mathcal E|\varepsilon} \arrow[r] \arrow[d] \& \unsure{\mathcal E|\varepsilon} \arrow[r] \arrow[d] \& 1 \\
1 \arrow[r] \& \unsure{T|\mathbb T}^I \arrow[r] \& S\Gamma \arrow[r] \arrow[d] \& \Gamma \arrow[r,no head] \arrow[d] \& 1 \\
{} \& {} \& 1 \& 1 \& {}
\end{tikzcd}}
so that we can consider $S\unsure{\mathcal E|\varepsilon}$ as an extension of $\Gamma$ by $\unsure{T|\mathbb T}^I\times\pi$
\[
1\longrightarrow \unsure{T|\mathbb T}^I\times\pi\longrightarrow S\unsure{\mathcal E|\varepsilon}\longrightarrow\Gamma\longrightarrow1
\]
from which \inserted{\unsure{the extension}} $\unsure{\mathcal E|\varepsilon}$ resp. $S\Gamma$ is obtained by dividing by $\unsure{T|\mathbb T}^{\unsure{I|S}}$ resp. by $\pi$.

\unsure{\deleted{Returning}|Returning} now to the model case \deleted{$U_{g,\nu}$}
$\unsure{U}_{g,\nu}=X_{g,\nu}\setminus S_{g,\nu}$, we \unsure{will use|use} \unsure{the|some} notations
\margin{\deleted{\unsure{$S\mathcal E_{g,\nu}$}}}
\[
\unsure{\mathcal E|\varepsilon}_{g,\nu}
=\mathrm{Aut}\deleted{\unsure{\mathrm{ext}}}_{\mathrm{lac}}
(\unsure{\pi|\Pi}_{g,\nu})
\]
($\stackrel{\mathrm{can}}{\simeq}\Gamma_{g,\nu+1}$ if $(g,\nu)\ne(0,0),(0,1)$ i.e. \deleted{\unsure{$\pi_{g,\nu}\ne0$}}) \illegible{1}

and
\[
S\Gamma_{g,\nu}=S\mathrm{Autext}_{\mathrm{lac}}(\unsure{\pi|\Pi}_{g,\nu})
\]
(extension of $\Gamma_{g,\nu}$ by $\unsure{T|\mathbb T}^I$)

and $S\unsure{\mathcal E|\varepsilon}_{g,\nu}$
($\simeq S\Gamma'_{g,\nu+1}/\unsure{T|\mathbb T|\Gamma}_{\unsure{s_\nu|g,\nu}}$
\inserted{\deleted{\illegible{1}} \unsure{where} $S\Gamma'_{g,\nu+1}$ denotes the}
subgroup of $S\Gamma'_{g,\nu+1}$ of the \unsure{elements} which fix \unsure{the} last element $s_\nu$ \ldots)